\section*{Bab \ref{ch:g5:areas} --- Keliling dan Luas}

\begin{solution}{exo:g5:areas:1}
Luas: $8 \times 5 = 40$ cm$^2$. \omterm{def:g4:measure:perimeter}{Keliling}:
$2 \times (8 + 5) = 26$ cm.
\end{solution}

\begin{solution}{exo:g5:areas:2}
$9 \times 9 = 81$ cm$^2$; \quad $12 \times 7 = 84$ m$^2$; \quad
$4.5 \times 6 = 27$ cm$^2$.
\end{solution}

\begin{solution}{exo:g5:areas:3}
Lebar: $63 \div 9 = 7$ cm. \omterm{def:g4:measure:perimeter}{Keliling}: $2 \times (9 + 7) = 32$ cm.
\end{solution}

\begin{solution}{exo:g5:areas:4}
Misalnya $6 \times 4$ (\omterm{def:g4:measure:perimeter}{kelilingnya} $20$) dan $8 \times 3$
(\omterm{def:g4:measure:perimeter}{kelilingnya} $22$): luasnya sama, $24$, \omterm{def:g4:measure:perimeter}{kelilingnya} berbeda.
\end{solution}

\begin{solution}{exo:g5:areas:5}
$3$ m$^2 = 30\,000$ cm$^2$; \quad $50\,000$ cm$^2 = 5$ m$^2$.
\end{solution}

\begin{solution}{exo:g5:areas:6}
Batang mendatar: $8 \times 2 = 16$ cm$^2$; batang tegak:
$2 \times 5 = 10$ cm$^2$; seluruhnya $26$ cm$^2$.
\end{solution}

\begin{solution}{exo:g5:areas:7}
$30 \times 20 - 24 \times 14 = 600 - 336 = 264$ cm$^2$ kayu.
\end{solution}

\begin{solution}{exo:g5:areas:8}
Setengah dari persegi panjang $8 \times 6$: $48 \div 2 = 24$ cm$^2$.
\end{solution}

\begin{solution}{exo:g5:areas:9}
Benih: luasnya $7 \times 4 = 28$ m$^2$, biayanya $28 \times 2 = 56$.
Pagar: \omterm{def:g4:measure:perimeter}{kelilingnya} $2 \times (7 + 4) = 22$ m, biayanya
$22 \times 3 = 66$.
\end{solution}

\begin{solution}{exo:g5:areas:10}
Dua ubin $50$ cm menjadi satu meter: sepanjang $12$ m ada $24$
ubin; sepanjang lebar $2$ m ada $4$ ubin. Seluruhnya:
$24 \times 4 = 96$ ubin.
\end{solution}

\begin{solution}{exo:g5:areas:11}
Sebelumnya: \omterm{def:g4:measure:perimeter}{kelilingnya} $2 \times (4 + 3) = 14$ cm, luasnya $12$
cm$^2$. Sesudahnya (sisi $8$ dan $6$): \omterm{def:g4:measure:perimeter}{kelilingnya} $28$ cm, luasnya
$48$ cm$^2$. \omterm{def:g4:measure:perimeter}{Kelilingnya} \emph{berlipat dua}; luasnya dikalikan
\emph{empat} ($2 \times 2$): panjang dan luas tidak berskala dengan
cara yang sama.

\end{solution}
